\chapter{Diseñando una DSL chica}
\label{cap:dsl-design}

\begin{intuicion}
La DSL de GraphHunter es deliberadamente \emph{chica}: el analista
debería poder escribir una hipótesis en menos de una línea sin
abrir manual. Eso es una decisión de diseño, no una limitación.
Una DSL chica gana legibilidad, reduce la curva de aprendizaje, y
fuerza al motor a ser bueno en pocos casos en lugar de mediocre en
muchos. Este capítulo explica los principios que guiaron el diseño
de la DSL: una sintaxis primaria de cadena, un solo separador (\texttt{;})
para extender a shapes graph, predicados como cláusulas
post-fix, y una semántica que mapea uno-a-uno a la estructura del
matcher.
\end{intuicion}

\section{Por qué una DSL en vez de SQL / Cypher / KQL}

GraphHunter podría haber expuesto su API en uno de los lenguajes
de query existentes:

\begin{itemize}
    \item \textbf{SQL}: estándar, pero verboso para multi-joins.
    Una hipótesis de 3 saltos = $\sim$15 líneas de SQL.
    \item \textbf{Cypher} (Neo4j): tiene sintaxis grafo
    (\texttt{(n)-[r]->(m)}), pero es enorme: una hipótesis simple
    obliga a aprender 30+ keywords, predicados completos, agrupaciones.
    \item \textbf{KQL} (Microsoft Sentinel): poderoso pero
    table-centric; modelar grafo encima de KQL pierde la naturaleza
    estructural.
\end{itemize}

\noindent
La decisión: una DSL nueva, optimizada para el problema concreto
(subgraph matching temporal), con una superficie chica que el
analista pueda dominar en una mañana.

\section{Principios de diseño}

\subsection{Principio 1: legibilidad de izquierda a derecha}

Una hipótesis se lee como se escribe: ``usuario que se autentica
en un host que se conecta a una IP''. La sintaxis refleja la
estructura del path:

\begin{lstlisting}
User -[Auth]-> Host -[Connect]-> IP
\end{lstlisting}

\noindent
Comparado con SQL (``\texttt{SELECT \dots FROM auth a JOIN connect c
ON a.host = c.host}''): el SQL invierte el orden mental ---primero
declara qué quiere, después cómo se conectan los datos---. La DSL
matchea el orden de la pregunta del analista.

\subsection{Principio 2: un solo separador}

La DSL legacy era pura cadena. La extensión \emph{shape graph} (L1
stage 3c) agregó la posibilidad de hipótesis no-cadena: estrellas,
árboles, cualquier hipergrafo \(\alpha\)-acíclico
(capítulo~\ref{cap:hipergrafos}). El diseño consideró varias
sintaxis alternativas:

\begin{enumerate}
    \item \texttt{User c \{ -[Auth]-> Host h1; -[Auth]-> Host h2 \}}:
    bloque explícito con corchetes.
    \item \texttt{User c -[Auth]-> Host h1 AND c -[Auth]-> Host h2}:
    \texttt{AND} entre cláusulas.
    \item \texttt{User c -[Auth]-> Host h1; c -[Auth]-> Host h2}:
    separador \texttt{;}.
\end{enumerate}

\noindent
Se eligió la 3 (separador \texttt{;}) por dos razones:

\begin{itemize}
    \item \textbf{Mínimo cambio sintáctico}. La cadena legacy NO
    contiene \texttt{;}, así que cualquier hipótesis con \texttt{;}
    es \emph{trivialmente} graph shape. Disambigua sin lookahead.
    \item \textbf{Familiaridad}. Los analistas que vienen de SQL
    asocian \texttt{;} con ``fin de declaración / siguiente
    cláusula''. El uso es coherente.
\end{itemize}

\subsection{Principio 3: predicados como cláusulas post-fix}

\begin{lstlisting}
User -[Auth]-> Host -[Connect]-> IP {k=2} HAVING count(distinct last) >= 5 WITHIN 30m
\end{lstlisting}

\noindent
La estructura del path va primero (la parte ``visual''); las
cláusulas modificadoras (\(k\)-simplicidad, ventana, agregaciones)
van al final, separadas por palabras clave (\texttt{WITHIN},
\texttt{HAVING}) o llaves (\texttt{\{k=N\}}).

Alternativa rechazada: predicados \emph{interlaced} con la cadena.
Hace el código más denso pero menos escaneable. La forma post-fix
se lee como un párrafo: ``este path; bajo estas condiciones''.

\subsection{Principio 4: defaults razonables}

Tres cosas que la DSL NO obliga a especificar:

\begin{itemize}
    \item \textbf{\(k\)-simplicidad}: default \(k = 1\) (path
    simple). Sólo \texttt{\{k=N\}} cuando el analista explícitamente
    quiere repeticiones.
    \item \textbf{Ventana temporal}: default \(\Delta t = +\infty\).
    Sólo \texttt{WITHIN} si el analista quiere acotar.
    \item \textbf{Aggregations}: default ``ninguna'', enumerar
    tal cual.
\end{itemize}

\noindent
La hipótesis mínima de 2 saltos ---\texttt{User -[Auth]-> Host
-[Connect]-> IP}--- es self-contained, sin cláusulas extras. Eso
baja la barrera de entrada: un analista nuevo puede escribir su
primera hipótesis sin leer manual.

\section{La gramática (informal)}

\begin{lstlisting}
hypothesis     := chain | graph
chain          := type_decl chain_step* trailing_clauses?
graph          := type_decl chain_step* (";" clause_head chain_step*)+ trailing_clauses?
chain_step     := "-[" rel_type ("{" predicates "}")? "]->" type_decl
type_decl      := EntityType (binding | "{" predicates "}")?
clause_head    := type_decl | binding_ref
binding_ref    := identifier   // referencia a binding ya declarado
predicates     := predicate ("," predicate)*
predicate      := key "=" value | key ">" value | ...
trailing_clauses := ("WITHIN" duration | "HAVING" agg | "{k=" int "}")*
\end{lstlisting}

\noindent
Tres terminales: \texttt{EntityType} (\texttt{User}, \texttt{Host},
etc.), \texttt{rel\_type} (\texttt{Auth}, \texttt{Connect}, etc.),
\texttt{identifier} (\texttt{c}, \texttt{h1}, \texttt{user}). El
parser dispatcha entre \texttt{type\_decl} y \texttt{binding\_ref}
en \texttt{clause\_head} consultando los bindings previamente
declarados (capítulo~\ref{cap:dsl-graph}).

\section{Lo que la DSL \emph{no} hace}

Decisiones de exclusión deliberadas:

\begin{itemize}
    \item \textbf{Sin disjunción}. No hay \texttt{OR} entre paths.
    Si el analista quiere ``A o B'', escribe dos hipótesis y une
    los resultados en el frontend. Razón: \texttt{OR} en DSL
    obliga al matcher a manejar disjunción, lo que rompe la
    semántica de path-único.
    \item \textbf{Sin negación}. No hay \texttt{NOT} sobre
    predicados (``vecino de tipo \emph{no}-IP''). Razón:
    negación general explota la cardinalidad ---el dominio de
    ``no-IP'' es enorme---. Caso de uso raro en threat hunting.
    \item \textbf{Sin variables libres en cláusulas trailing}. El
    \texttt{HAVING} sólo refiere a \texttt{first}, \texttt{last},
    o agregaciones builtin. No hay binding de variables
    intermedias. Razón: simpleza del parser y evaluador.
    \item \textbf{Sin recursividad}. No hay
    \texttt{User -[Knows]*-> User} (transitivo de longitud
    variable). Razón: explosión de estado; los hunts realmente
    útiles son de longitud chica.
\end{itemize}

\noindent
Cada exclusión es un trade-off: alguna pregunta legítima no se
puede expresar. La regla: si el caso de uso aparece en \(<\) 5\%
de los hunts del corpus, no entra en la DSL; el analista lo resuelve
combinando varias queries.

\section{La sintaxis no es la semántica}

Importante: la DSL es la \emph{interfaz}, no la
\emph{representación interna}. El parser convierte la string a
\texttt{Hypothesis} (capítulo~\ref{cap:ast}); a partir de ahí
\emph{todo} el motor trabaja sobre \texttt{Hypothesis}, no sobre
strings. Eso permite:

\begin{itemize}
    \item Construir hipótesis programáticamente desde otros
    transports (MCP, REST) sin parsear DSL.
    \item Cambiar la sintaxis de la DSL sin tocar el matcher.
    \item Generar la DSL desde un constructor visual (planeado
    futuro UI).
\end{itemize}

\section{Resumen}

\begin{itemize}
    \item DSL chica deliberada: legible, low-barrier, mapea a la
    estructura del matcher.
    \item Cuatro principios: legibilidad LR, un solo separador,
    predicados post-fix, defaults razonables.
    \item Exclusiones deliberadas: sin disjunción, sin negación,
    sin recursividad, sin variables libres en trailing.
    \item DSL es la interfaz; \texttt{Hypothesis} es la
    representación interna. Cualquier transport puede saltarse el
    parser.
\end{itemize}

\begin{ejercicios}
\begin{enumerate}
    \item \starred{} Escribí en la DSL la hipótesis: ``un proceso
    que escribió un archivo y después se conectó a una IP, todo
    en menos de 5 minutos''. Identificá los tres componentes
    (cadena, ventana, predicados).

    \item ¿Por qué \texttt{;} es mejor disambiguador que
    \texttt{AND}? Discutí dos casos donde \texttt{AND} sería
    ambiguo.

    \item \starred{} Diseñá una sintaxis para
    \emph{disjunción} (``patrón A o patrón B''). ¿Cómo se
    parsearía? ¿Qué cambia para el matcher?

    \item Imaginá un caso de uso real (de threat hunting) que la
    DSL actual \emph{no} puede expresar. ¿Cuál es? ¿Vale la pena
    extender la DSL para cubrirlo?

    \item Si el frontend quisiera ofrecer un constructor visual
    (drag-and-drop nodes y edges) en vez de DSL, ¿qué tendría que
    generar para que el motor lo entienda?
\end{enumerate}
\end{ejercicios}
